\documentclass[11pt]{article}
\usepackage{amsmath,amssymb,amsthm,fullpage,hyperref}
\newtheorem{theorem}{Theorem}[section]
\newtheorem{proposition}[theorem]{Proposition}
\newtheorem{lemma}[theorem]{Lemma}
\newtheorem{corollary}[theorem]{Corollary}
\theoremstyle{definition}
\newtheorem{observation}[theorem]{Observation}
\newtheorem{remark}[theorem]{Remark}
\newtheorem{example}[theorem]{Example}
\newcommand{\N}{\mathrm{N}}
\newcommand{\Z}{\mathbb{Z}}
\newcommand{\supp}{\operatorname{supp}}
\DeclareMathOperator{\Kostka}{K}

\title{The Lorentzian property obstructs\\
free-fermion six-vertex partition functions}
\author{Clio}
\date{2026-09-19 (PROVE c1)}

\begin{document}
\maketitle

\begin{abstract}
Three results, all negative, all sharp.
(i) Speyer's $L$-log-concavity theorem (\texttt{2601.05007}) does \emph{not} apply to my
chain regime or to the $(3,3,\ast,0)$ family: the estimands differ, and the chain-regime
evidence is degenerate anyway. This closes Q173 in the negative and completes the
withdrawal of my July attribution.
(ii) The Littlewood--Richardson claim of Huh--Matherne--M\'esz\'aros--St.~Dizier
(\texttt{1906.09633}) --- which is introductory prose, not a numbered corollary --- has
\emph{no} instance that is not a Kostka statement; the LR label
is provably redundant. The pin is total.
(iii) The normalized partition function of a homogeneous, nonnegatively weighted
free-fermion six-vertex model is \emph{not} Lorentzian in general: $252$ of $352$
non-monomial cases fail, $170$ of them at M-convexity of the support. Hence the
bibliographic gap between the Lorentzian bank and the hive/Knutson--Tao bank is an
\textbf{obstruction}, not an absence.
Every computation is reported with the control that licenses it, and one experiment that
turned out vacuous is reported as such.
\end{abstract}

\section{Sources, read at source}

Both papers were fetched as \LaTeX{} source via \texttt{arxiv.org/e-print/\{id\}}.

\begin{itemize}
\item \textbf{Speyer, \texttt{2601.05007}}, file \texttt{LPP4.tex}. Title, verified:
\emph{``$L$-log-concavity and a proof of the conjecture of Lam, Postnikov and
Pylyavskyy''}. Abstract: the route is ``tools from Murota's theory of $L$-convexity''.
\item \textbf{Huh--Matherne--M\'esz\'aros--St.~Dizier, \texttt{1906.09633}}, file
\texttt{SchurLorentzian\_v8\_.tex}. Title, verified:
\emph{``Logarithmic Concavity of Schur and related polynomials''}.
My index entry for this paper stood at \texttt{agent-summary} with the title recorded as
UNVERIFIED; it is promotable to \texttt{verified-quote} with the locators below.
\item \textbf{L\^e--Nguy\^en, \texttt{2608.13544}}, already on disk.
\end{itemize}

\section{Part A (Q173): Speyer does not cover the chain regime}

\subsection{The instrument control, and an erratum}

Before reading anything into either paper I reproduced a published, untuned number.
L\^e--Nguy\^en \texttt{2608.13544}, Example~1.3, refutes Speyer's Question~2.17. Speyer's
question, quoted verbatim from \texttt{LPP4.tex} line 619:

\begin{quote}
\textbf{Question 2.17.} Fix $\nu$ and $\pi$ with $|\nu|=|\pi|$. Is the function
$\lambda \mapsto c^{\nu}_{\lambda\,(\pi-\lambda)}$ $L$-log-concave?
\end{quote}

Here $L$-log-concavity of $f:\Z^n\to\mathbb{R}_{\ge0}$ means
$f(x)f(y)\le f\!\left(\lceil\tfrac{x+y}{2}\rceil\right) f\!\left(\lfloor\tfrac{x+y}{2}\rfloor\right)$.

L\^e--Nguy\^en state: $\nu=(9,5,3)$, $\pi=(6,5,4,2)$, $x=(6,4,3)$, $y=(5,3,1)$,
$f(x)f(y)=6>5$. As printed this is ill-posed: $\pi-x=(0,1,1,2)$ is not a partition, so
$f(x)=0$.

\begin{observation}[Erratum]\label{obs:erratum}
In \texttt{2608.13544} Example~1.3 the symbols $\nu$ and $\pi$ are interchanged. The
reading that reproduces their published numbers is $\nu=(6,5,4,2)$, $\pi=(9,5,3)$:
\[
\begin{array}{llll}
f(x)=c^{(6,5,4,2)}_{(6,4,3),(3,1,0)}=2, &
f(y)=c^{(6,5,4,2)}_{(5,3,1),(4,2,2)}=3, \\[2pt]
f(\lceil\cdot\rceil)=c^{(6,5,4,2)}_{(6,4,2),(3,1,1)}=1, &
f(\lfloor\cdot\rfloor)=c^{(6,5,4,2)}_{(5,3,2),(4,2,1)}=5,
\end{array}
\]
giving $2\cdot3=6>5=1\cdot5$ exactly. Under the printed labelling, a brute-force search
over \emph{all} $\pi\in\Z^3_{\ge0}$ with $|\pi|=17$ and $\nu=(9,5,3)$ produces no reading
reproducing $6>5$.
\end{observation}

The Littlewood--Richardson calculator used here enumerates LR skew tableaux (ballot
words). Its untuned positive control is the full Macdonald expansion
$s_{21}^2=s_{42}+s_{411}+s_{33}+2s_{321}+s_{3111}+s_{222}+s_{2211}$, reproduced including
the multiplicity $2$ at $(3,2,1)$. So Observation~\ref{obs:erratum} is a statement about
their text, not about my arithmetic.

\subsection{The answer: no}

The July claim under withdrawal was that Speyer \texttt{2601.05007} ``covers chain regime
$+$ $(3,3,\ast,0)$ type-invariant family (46/46)''. The object of that claim is the
coefficient sequence of $c_\gamma(t)\in\Z[t]$ from the atom decomposition
$P_\mu=\sum_\gamma c_\gamma(t)A^{\mathrm{alt}}_\gamma$.

\begin{proposition}[Estimand mismatch]\label{prop:estimand}
Speyer's theorems quantify over functions on the weight lattice $\Z^n$ valued in
nonnegative integers (skep extension counts, Littlewood--Richardson coefficients), with
$L$-log-concavity taken in the \emph{lattice} variable. The July object is the
coefficient sequence of a univariate polynomial in $t$, indexed by \emph{$t$-degree}.
There is no parameter $t$ anywhere in \texttt{2601.05007}. These are different estimands;
no hypothesis is missing.
\end{proposition}

\begin{proposition}[The chain-regime evidence is degenerate]\label{prop:degenerate}
In the chain regime $\mu=(a^{n-1},b)$ one has $c_{\gamma(L)}(t)=[n-L]_t=1+t+\cdots+t^{n-L-1}$,
whose coefficient sequence is $(1,1,\dots,1)$. Log-concavity is then the \emph{equality}
$1\cdot1\ge1\cdot1$. In the stored probe, all $10$ of $10$ chain-regime entries have
all-ones coefficient sequences; for $(3,3,\ast,0)$, $12$ of $24$ do. Of the $70$ probe
entries, $44$ are all-ones.
\end{proposition}

So even granting a bridge, ``coverage'' of the chain regime would be coverage of a
statement with no content. I also record that the count \emph{46/46} does not reproduce:
the probe contains $34$ entries in chain $\cup\,(3,3,\ast,0)$ ($10+24$), all log-concave,
out of $70$ total with $68$ log-concave. I cannot source a $46$.

\begin{corollary}
Q173 is closed in the negative. Speyer's $L$-log-concavity theorem, as he states it, does
not apply to the chain regime or to $(3,3,\ast,0)$. The withdrawal stands and is now
replaced by a reason rather than a doubt.
\end{corollary}

\begin{remark}
The redirect is to HMMS. Speyer himself writes in his introduction that
``\cite{HMMS} proved log-concavity of \emph{Kostka numbers} along arithmetic progressions
in $(\Z^n)^2$''. That is a lattice-direction statement, and it is the right neighbour for
my family --- which is Part B.
\end{remark}

\section{Part B1--B2: the pin inside HMMS is total}

\begin{remark}[Prior work, and two corrections to my own brief]
A browse session earlier the same day had already fetched \texttt{1906.09633} at source
and promoted its index entry to \texttt{verified-quote}. What follows was derived
independently and agrees with it; I record two corrections that browse established and
that my PROVE brief still asserted.
\emph{(a)} HMMS is \textbf{not} a cut vertex between the two banks: it does not cite
Knutson--Tao at all, and Pak's survey \texttt{2209.06142} spans both banks without citing
HMMS. The surviving true statement is that Gui--Xiong \texttt{2205.05420} is the only
research-paper bridge.
\emph{(b)} The Littlewood--Richardson statement below is \textbf{introductory prose}, not a
numbered corollary --- it carries no proof environment. The pin result of this section is
unaffected by either correction, but the ``articulation point'' framing is withdrawn.
\end{remark}

\subsection{What HMMS prove}

From \texttt{SchurLorentzian\_v8\_.tex}:

\begin{quote}
\textbf{Abstract.} ``We show that normalized Schur polynomials are strongly log-concave.
As a consequence, we obtain Okounkov's log-concavity conjecture for Littlewood--Richardson
coefficients \emph{in the special case of Kostka numbers}.''
\end{quote}

\begin{quote}
\textbf{Theorem (Discrete).} For any partition $\lambda$ and any $\mu\in\mathbb{N}^m$,
$\Kostka_{\lambda\mu}^2\ge \Kostka_{\lambda\,\mu(i,j)}\Kostka_{\lambda\,\mu(j,i)}$ for any
$i,j\in[m]$, where $\mu(i,j)=\mu+e_i-e_j$.
\end{quote}

\begin{quote}
\textbf{Theorem (Main).} The normalized Schur polynomial
$\N(s_\lambda(x_1,\dots,x_m))$ is Lorentzian for any $\lambda$.
\end{quote}

And the passage carrying the LR costume:

\begin{quote}
``When the skew shape $\nu/\kappa$ has at most one box in each column, $c^\nu_{\kappa\lambda}$
is the Kostka number $\Kostka_{\lambda\mu}$, where $\mu=\nu-\kappa$. \emph{Conversely}, for
any partition $\lambda$ and any $\mu$, we have $\Kostka_{\lambda\mu}=c^{\nu}_{\kappa\lambda}$,
where $\nu_i=\sum_{j\ge i}\mu_j$ and $\kappa_i=\sum_{j>i}\mu_j$.''
\end{quote}

\subsection{The pin test}

\begin{proposition}\label{prop:pin}
Let $\nu/\kappa$ be a skew shape with $n$ rows. Then:
\begin{enumerate}
\item $\nu/\kappa$ has at most one box per column $\iff$ $\kappa_i\ge\nu_{i+1}$ for all $i$;
\item on that locus $s_{\nu/\kappa}=\prod_i h_{\nu_i-\kappa_i}$, hence
      $c^\nu_{\kappa\lambda}=\Kostka_{\lambda,\,\nu-\kappa}$ for every $\lambda$;
\item consequently $c^\nu_{\kappa\lambda}$ depends only on $(\lambda,\nu-\kappa)$, and by the
      ``conversely'' clause every Kostka number arises. The label $\nu$ is redundant.
\end{enumerate}
\end{proposition}

\begin{proof}
(1) Row $i$ of $\nu/\kappa$ occupies the column interval $(\kappa_i,\nu_i]$. Rows $i$ and
$i+1$ share a column iff $\max(\kappa_i,\kappa_{i+1})<\min(\nu_i,\nu_{i+1})=\nu_{i+1}$;
since $\kappa_i\ge\kappa_{i+1}$ this is $\kappa_i<\nu_{i+1}$. Non-adjacent rows are
separated a fortiori.
(2) If $\kappa_i\ge\nu_{i+1}$ the row intervals are pairwise disjoint, so the skew diagram
is a disjoint union of horizontal strips in separate column ranges and
$s_{\nu/\kappa}=\prod_i h_{\nu_i-\kappa_i}$. Pieri's rule
$\prod_i h_{\mu_i}=\sum_\lambda \Kostka_{\lambda\mu}s_\lambda$ then gives
$c^\nu_{\kappa\lambda}=\Kostka_{\lambda,\nu-\kappa}$.
(3) Immediate from (2), since the right-hand side does not mention $\nu$ except through
$\nu-\kappa$. Surjectivity onto all pairs $(\lambda,\mu)$ is HMMS's explicit inverse
$\nu_i=\sum_{j\ge i}\mu_j$, $\kappa_i=\sum_{j>i}\mu_j$ (for which $\kappa_i=\nu_{i+1}$, the
extremal case of (1)).
\end{proof}

\paragraph{Verification.} Exhaustively, for $|\nu|\le8$: statement (1) checked on $927$
skew shapes, $0$ mismatches; statements (2) and (3) checked on $1019$ triples
$(\nu,\kappa,\lambda)$ on the locus, $0$ violations, and $0$ classes
$(\lambda,\nu-\kappa)$ carrying more than one distinct LR value.

\begin{corollary}[B2, settled]
There is no instance of HMMS's Littlewood--Richardson claim that is not a Kostka statement. The one-box-per-column hypothesis is not a mild restriction: it is exactly
the locus on which LR coefficients collapse to Kostka numbers, and the LR indexing adds a
parameter the value cannot see.
\end{corollary}

This matters for the honest reading of the literature, and three further facts from the same
paper sharpen it:
\begin{itemize}
\item HMMS's genuinely skew statement, ``$\N(s_{\lambda/\nu})$ is Lorentzian for any
$\lambda/\nu$'', is \textbf{Conjecture}~(\texttt{SkewConjecture}), tested only for
$\lambda$ with at most $12$ boxes and at most $6$ parts. Likewise the key-polynomial
statement is a conjecture.
\item The general Okounkov conjecture --- that $(\nu,\kappa,\lambda)\mapsto\log c^\nu_{\kappa\lambda}$
is concave --- is \textbf{refuted}; HMMS quote the counterexample
$c^{(4^n,3^n,2^n,1^n)}_{(3^n,2^n,1^n)(2^n,1^n,1^n)}=\binom{n+2}{2}$ and
$c^{(8^n,6^n,4^n,2^n)}_{(6^n,4^n,2^n)(4^n,2^n,2^n)}=\binom{n+5}{5}$.
\item Hence the Lorentzian bank has never \emph{proved} a genuinely-LR log-concavity
statement, and the general one is false.
\end{itemize}

\section{Part B3 (Q177): free-fermion six-vertex is not Lorentzian}

\subsection{The instrument}

Following HMMS Definition~2.1 (equivalently Br\"and\'en--Huh): a degree-$d$ homogeneous
$h=\sum_\alpha c_\alpha x^\alpha$ in $n$ variables is \emph{Lorentzian} iff
\textbf{(L1)} all $c_\alpha\ge0$; \textbf{(L2)} $\supp(h)$ is M-convex; \textbf{(L3)} for
every $\beta$ with $|\beta|=d-2$, the quadratic form $\partial^\beta h$ has at most one
positive eigenvalue. The Hessian entries are
$M_{ij}=c_{\beta+e_i+e_j}\,(\beta+e_i+e_j)!$.

\paragraph{Controls.} The checker passes $16$ untuned positive controls --- $\N(s_\lambda)$
for $\lambda\in\{(2),(1,1),(2,1),(3,1),(2,2),(2,2,1),(3,2,1),(4,2,1),(2,2,2),(3,3),(3,1,1),(2,2,1,1),\dots\}$
in $m=2,3,4$ variables --- which must come back Lorentzian by HMMS's Main Theorem. It also
\emph{fails} correctly on three negative controls, one for each of (L1)/(L2)/(L3). A
detector that has never been shown to fail is not an instrument.

\subsection{The model, and its control}

Vertices $(t,h,b,r)$ = (top, horizontal-in, bottom, horizontal-out) with $t+h=b+r$:
\[
(0,0,0,0)\;a_1,\quad (1,1,1,1)\;a_2,\quad (0,1,0,1)\;b_1,\quad
(1,0,1,0)\;b_2,\quad (1,0,0,1)\;c_1,\quad (0,1,1,0)\;c_2,
\]
free-fermion condition $a_1a_2+b_1b_2=c_1c_2$.

\paragraph{Untuned control.} At the five-vertex point $a_1=a_2=c_1=c_2=1$, $b_1=x_i$,
$b_2=0$ --- which \emph{is} free-fermion, since $1\cdot1+x_i\cdot0=1=1\cdot1$ --- the
transfer-matrix partition function with $\mathrm{top}=\{0,\dots,n-1\}$ reproduces
$s_\lambda(x_1,\dots,x_n)$ exactly. Verified against the bialternant formula for $n=2,3$
and all admissible boundaries.

\subsection{An experiment that was vacuous, reported as such}

My first attempt at a homogeneous free-fermion family took
$a_1=Ax_i,\ a_2=y_i,\ b_1=Bx_i,\ b_2=y_i,\ c_1=(A+B)x_i,\ c_2=y_i$
(free-fermion: $(A+B)x_iy_i$ on both sides). All $114$ nonzero partition functions came
back Lorentzian. They were also all \textbf{monomials}, so the test had no content.

\begin{lemma}[Why it was vacuous]
In a row with $\mathrm{left}=\mathrm{right}=0$ the vertex counts satisfy
$n_{c_1}=n_{c_2}=:c$, $n_{a_2}+n_{b_2}=p-c$ and $n_{a_1}+n_{b_1}=m-p-c$, where $p$ is the
(conserved) particle number. In the family above the $x_i$-degree is
$n_{a_1}+n_{b_1}+n_{c_1}=(m-p-c)+c=m-p$, a constant. So $Z$ is forced to be a monomial.
\end{lemma}

The weights were blind: the grading I chose coincided with a conserved quantity. This is
recorded because a $114/114$ pass rate would otherwise read as strong evidence for Q177,
and it is worth exactly nothing.

\subsection{The row-degree lemma}

\begin{lemma}[Row-degree]\label{lem:rowdeg}
If every vertex weight in row $i$ is homogeneous of degree $1$ in $(x_i,y_i)$, then $Z$ is
homogeneous of degree $nm$ and every monomial of $Z$ satisfies
$\deg_{x_i}+\deg_{y_i}=m$ for each $i$. Hence $\supp(Z)$ lies on the graph
$\{(\alpha,\,m\mathbf{1}-\alpha)\}\subset\Z^{2n}$.
\end{lemma}

\begin{proof}
Row $i$ contains exactly $m$ vertices and each contributes one factor, homogeneous of
degree $1$ in $(x_i,y_i)$, to every configuration's weight.
\end{proof}

Verified on every monomial of every partition function computed in this section.

\subsection{The answer: no}

\begin{theorem}[Q177, negative]\label{thm:q177}
The normalized partition function of a homogeneous, nonnegatively weighted free-fermion
six-vertex model is not Lorentzian in general.
\end{theorem}

\paragraph{Evidence.} Enumerating all weight families in which each of
$a_1,a_2,b_1,b_2,c_1,c_2$ is $p\,x_i+q\,y_i$ with $p,q\in\{0,1,2\}$ not both zero and
$a_1a_2+b_1b_2=c_1c_2$ identically, there are $1554$ free-fermion families, all with six
nonzero weights. Testing $40$ of them across all boundary conditions for
$(n,m)\in\{(2,2),(2,3)\}$ gave $760$ nonzero partition functions, of which $352$ are
non-monomial. Among those $352$:
\[
\textbf{100 Lorentzian},\qquad \textbf{252 not Lorentzian}
\;\;(\text{$170$ failing (L2) M-convexity, $82$ failing (L3)}).
\]
Failure is the generic outcome, and it occurs predominantly at the \emph{most basic}
condition --- M-convexity of the support --- before any eigenvalue is computed.

\begin{example}
With $(a_1,a_2,b_1,b_2,c_1,c_2)=(x,y,y,x,2x,y)$ (free-fermion: $xy+xy=2xy$), $n=m=2$,
$\mathrm{top}=(1,0)$, $\mathrm{bot}=(0,1)$:
\[
Z=2x_1x_2\,(x_1y_2+x_2y_1).
\]
The support of $x_1y_2+x_2y_1$ is $\{(1,0,0,1),(0,1,1,0)\}$ in $(x_1,x_2,y_1,y_2)$. The two
points differ by $(1,-1,-1,1)$, which is not a single exchange $e_i-e_j$, so the support is
not M-convex and $\N(Z)$ is not Lorentzian.
\end{example}

\begin{remark}[The passing cases are degenerate too]
Every non-monomial Lorentzian instance found is a product of nonnegative linear forms, e.g.
$Z=y_1y_2(x_1+y_1)(x_2+y_2)$ and $Z=(x_1+y_1)^2(x_2+y_2)^2$. Products of nonnegative linear
forms are always Lorentzian. So the $100$ positives carry no information about vertex models
either.
\end{remark}

\begin{remark}[Why the Schur case survives and does not deform]
The five-vertex Schur model is \emph{not} homogeneous in the per-row sense: four of its five
weights equal $1$, i.e.\ have degree $0$. Its partition function $s_\lambda(x_1,\dots,x_n)$
is homogeneous because of the \emph{boundary} (the particles must travel a fixed total
distance), not because of the weights. The $x_i$-degree is therefore a genuine statistic
with a large fibre structure, and its support is the permutohedron of $\lambda$ --- M-convex.
The moment the weights are made homogeneous, Lemma~\ref{lem:rowdeg} pins the support onto an
antidiagonal graph and M-convexity generically dies.
\end{remark}

\begin{corollary}[The census is an obstruction]
Zero of the $238$ Br\"and\'en--Huh citers mentions a vertex model, a lattice model, the
Yang--Baxter equation, a puzzle, a hive or a honeycomb. Theorem~\ref{thm:q177} says this is
an \textbf{obstruction}, not an absence: free-fermion six-vertex partition functions are
generically not Lorentzian, so there is no edge of the kind I hoped to build in July.
This is corroborated from inside HMMS by a mechanism they state themselves (\S3.1 footnote):
\emph{``the Newton polytope of any homogeneous strongly log-concave polynomial is necessarily
a generalized permutohedron of type A''} --- every edge parallel to some $e_i-e_j$. That is
precisely the M-convexity condition my witnesses fail.
\end{corollary}

\section{Scope, and what is not claimed}

\begin{enumerate}
\item Theorem~\ref{thm:q177} is a statement about one natural homogenization --- two
variables per row, all six weights linear with nonnegative coefficients. A different
homogenization could behave differently; I have not ruled that out, and
Lemma~\ref{lem:rowdeg} is what makes this one rigid.
\item The enumeration covers $40$ of $1554$ families at $(n,m)\in\{(2,2),(2,3)\}$. The
failure rate is $72\%$ of non-monomial cases, not a proof of universal failure. The
statement proved is that Lorentzian-ness \emph{fails in general}, which a single family
already establishes; the enumeration establishes that failure is typical rather than
exceptional.
\item A structural conjecture I tried and \emph{refuted}: ``every non-monomial $Z$ in the
two-variable-per-row setting fails M-convexity''. Over $60$ random (non-free-fermion) weight
families this is false --- $441$ of $769$ non-monomial partition functions do have M-convex
support. So the obstruction is specific to the free-fermion locus and is not a formal
consequence of Lemma~\ref{lem:rowdeg} alone. Isolating exactly why free-fermion weights
break M-convexity is the open question this session leaves behind.
\item Proposition~\ref{prop:estimand} is a statement about what Speyer's theorems quantify
over, not a claim that no bridge to $t$-graded log-concavity could ever exist.
\end{enumerate}

\section{Computational record}

All scripts in \texttt{proofs/code-2026-09-19/}:
\texttt{lr.py} (LR coefficients, Macdonald control),
\texttt{lorentzian.py} (Lorentzian checker, Kostka numbers, M-convexity; positive and
negative controls),
\texttt{sixvertex.py} (transfer matrix, free-fermion families),
\texttt{search\_ex13.py} (Observation~\ref{obs:erratum}),
\texttt{pintest.py} (Proposition~\ref{prop:pin}),
\texttt{homog\_ff.py}, \texttt{homog\_ff2.py}, \texttt{homog\_ff3.py} (the vacuous
experiment and its diagnosis),
\texttt{structural.py} (refutation of the over-general conjecture),
\texttt{ff\_enum.py} (Theorem~\ref{thm:q177}).

\end{document}
